\section{Causal Discovery and Time Series: Progress and Key Findings}
\label{sec:survey}

Causal discovery has developed from classical graphical models into methods that
handle mixed data, heterogeneity, and temporal dependence. Recent work focuses on
nonstationary and high-dimensional time series.

\subsection{Graphical Causal Discovery}

The main classical approaches are constraint-based methods, such as PC and FCI,
and score-based DAG learning. These methods provide the theoretical basis for
structure learning, but their reliability can decrease with nonlinearity,
measurement error, nonstationarity, and latent confounding
\cite{glymour2019review}. Hybrid approaches combine the strengths of both families.
For example, constraint-based filtering can reduce the search space before Bayesian
sampling is applied \cite{kuipers2018efficient}.

Recent methods also address mixed data and larger graphs. Likelihood-ratio tests
improve conditional-independence testing for mixed variables
\cite{tsagris2018constraint}, while variational Bayesian methods improve scalability
\cite{annadani2021variational}. Other work studies assumption violations
\cite{faltenbacher2025internal}, cyclic graphs \cite{mooij2020fci}, joint graph and
effect uncertainty \cite{castelletti2021structural}, and NOTEARS-style learning for
mixed data \cite{zhao2024notearsm}.

\subsection{Heterogeneity, Non-stationarity, and Mechanism Shifts}

Changes across environments can provide information about causal direction. CD-NOD
uses distribution shifts to recover causal structure from heterogeneous or
nonstationary data \cite{huang2020causal}. The Sparse Mechanism Shift hypothesis
shows that a full graph can become identifiable when only a small number of causal
mechanisms change across environments \cite{perry2022causal}.

Recent methods also cluster observations by their causal mechanisms. HCL learns
latent clusters and their graphs from mixed observational data without environment
labels \cite{li2025hcl}. MCVCI and MCVCC extend this idea to mixtures of bivariate
mechanisms \cite{liu2025hybrid}. These methods motivate the present work, which
considers unknown regime counts and sequential observations.

\subsection{Causal Discovery for Time Series (Stationary and Non-stationary)}

Time-series causal discovery includes Granger-based, constraint-based, noise-based,
score-based, logic-based, topology-based, and difference-based methods
\cite{assaad2022survey}. No family is best in every setting. Each depends on
assumptions such as linearity, causal sufficiency, stationarity, or acyclicity.
Multivariate time series add further problems, including subsampling and unmeasured
confounding \cite{glymour2019review}.

CD-NOTS extends CD-NOD to nonstationary time series
\cite{sadeghi2025causal,huang2020causal}. It models contemporaneous and lagged
effects while enforcing consistency across time. It also corrects problems that can
arise when lagged variables are used in conditional-independence tests. The method
supports KCIT, RCoT, CMIknn, and ParCorr, allowing the test to be selected for the
sample size and dimensionality. Its financial-data experiments show that
nonstationary and nonlinear effects can be missed by stationary linear methods.

\subsection{Deep-learning and Granger-based Methods}
\label{ssec:granger}

Deep-learning methods extend Granger-based discovery to larger and more complex
time series. TCDF uses an attention-based temporal CNN and permutation validation
to learn graphs and lags \cite{nauta2019causal}. NTiCD combines LSTMs with graph
convolutions \cite{absar2023neural}. CUTS+ uses Granger-based neural components
with graph neural networks for irregular streams \cite{cheng2023cuts}. Transformer
attention has also been used to represent changing Granger graphs
\cite{zhu2025attention}. These approaches depend on temporal ordering and therefore
do not directly apply to unordered i.i.d. observations.

\subsection{Temporal Causal Discovery with Unknown Regimes: CASTOR}

CASTOR is the closest existing method to this thesis. It learns causal graphs from
multivariate time series with an unknown number of sequential regimes
\cite{rahmani2025castor}. It estimates the regime count, the boundaries, and the
causal DAG for each regime.

CASTOR uses expectation-maximisation. The E-step updates regime-membership
probabilities, while the M-step estimates a BIC-penalised DAG under the NOTEARS
acyclicity constraint \cite{zheng2018dags}. It starts with equal-width windows and
merges adjacent windows with similar causal structures. The method provides
identifiability results under Gaussian equal-variance noise. Experiments on graphs
with up to 40 nodes and two real benchmarks report improvements over CD-NOD,
RPCMCI, and DYNOTEARS in regime recovery and graph recovery.

CASTOR and LIRA address the same broad problem: finding an unknown number of
sequential regimes and their graphs without environment labels. Their boundary
criteria differ. CASTOR accepts boundaries through a global EM and BIC objective.
LIRA first applies a Chow $F$-test to ask whether the coefficient vectors differ,
then uses BIC to control model complexity. This gives a statistical test and a
separate parsimony check for each candidate boundary.

LIRA also uses a domain-specific structural model derived from the Fundamental
Plane of Black Hole Activity. This fixes the variable set and gives the model
parameters a physical interpretation. CASTOR is domain-agnostic and therefore more
general. In return, LIRA uses stronger physical assumptions. Its greedy binary
segmentation may miss jointly detectable boundaries, and its piecewise-constant
model cannot represent gradual changes. These limitations motivate future
comparisons with CASTOR and DyCAST.

\subsection{Continuously Evolving Causal Structure: DyCAST}

DyCAST models causal structure as a continuous function of time rather than as a
sequence of fixed regimes \cite{cheng2025dycast}. Neural ODEs constrained to the
DAG manifold allow it to represent smooth transitions. Experiments on synthetic,
fMRI, CausalTime, and motion-capture data show strong performance and better
scaling than several static temporal methods. However, DyCAST also requires
temporal ordering and therefore is not a direct solution for unordered data.

\subsection{Challenges and Methodological Themes}

The performance of causal discovery depends strongly on whether the model
assumptions match the data \cite{assaad2022survey}. Important challenges include
nonlinearity, causal insufficiency, mixed data types, high dimensionality, and
finite-sample error. Existing responses include hybrid Bayesian learning
\cite{kuipers2018efficient}, mixed-data independence tests
\cite{tsagris2018constraint}, and graph neural networks for coarse-to-fine discovery
\cite{cheng2023cuts}.

The field is also moving toward uncertainty quantification. Variational causal
networks estimate distributions over possible DAGs \cite{annadani2021variational},
and joint graph-effect models propagate structural uncertainty into effect estimates
\cite{castelletti2021structural}. Validation is another growing focus, with evidence
that learned Granger graphs can improve forecasting \cite{zhang2023causalformer}.

\subsection{Conclusion of the Survey}

Causal discovery now includes methods for mixed data, heterogeneity, uncertainty,
and temporal dynamics. However, each method depends on particular assumptions.
Many approaches require environment labels, an auxiliary time or domain index, or
temporal ordering. This thesis focuses on discovering causal regimes from
unlabelled observations while estimating the number of regimes and testing whether
their mechanisms differ. CASTOR provides the closest temporal analogue, but its
dependence on sequential regime assignment does not solve this problem directly.

The main related methods are compared in Table~\ref{tab:related_works}.
\newcolumntype{P}[1]{>{\raggedright\arraybackslash}p{#1}}

{\small
\renewcommand{\arraystretch}{1.3}
\begin{longtable}{|P{2.0cm}|P{2.3cm}|P{2.5cm}|P{1.8cm}|P{4.3cm}|}
\caption{Summary of major related works in causal discovery, regime detection,
and causal clustering.}
\label{tab:related_works} \\
\hline
\textbf{Study} & \textbf{Task} & \textbf{Method} & \textbf{Setting} & \textbf{Main Finding / Gap} \\
\hline
\endfirsthead

\multicolumn{5}{c}%
{{\bfseries \tablename\ \thetable{} -- continued from previous page}} \\
\hline
\textbf{Study} & \textbf{Task} & \textbf{Method} & \textbf{Setting} & \textbf{Main Finding / Gap} \\
\hline
\endhead

\hline \multicolumn{5}{|r|}{{Continued on next page}} \\ \hline
\endfoot

\hline
\endlastfoot

Spirtes et al.\ \cite{spirtes2000causation}
& DAG learning from observational data
& PC and FCI constraint-based algorithms
& i.i.d., single environment
& Foundational framework for causal discovery. Assumes a single fixed causal model throughout. \\
\hline
Zheng et al.\ \cite{zheng2018dags}
& Continuous DAG structure learning
& NOTEARS: smooth algebraic acyclicity constraint
& i.i.d., single environment
& Enables gradient-based DAG learning at scale. The acyclicity constraint underpins DYNOTEARS and CASTOR. \\
\hline
Pamfil et al.\ \cite{pamfil2020dynotears}
& Time-series causal graph recovery
& DYNOTEARS: NOTEARS extended to contemporaneous and lagged edges
& Stationary time series
& Efficiently recovers window causal graphs. Assumes stationarity holds throughout the entire series. \\
\hline
Huang et al.\ \cite{huang2020causal}
& Causal discovery under heterogeneity
& CD-NOD: nonparametric method exploiting distribution shifts
& Multiple domains or time-indexed data
& Heterogeneity aids causal identification. Requires an auxiliary domain index or time stamp. \\
\hline
Peters et al.\ \cite{peters2016causal}
& Causal parent identification
& ICP: invariant prediction across environments
& Multiple labelled environments
& Exploits causal invariance effectively. Requires environment labels as a strict prerequisite. \\
\hline
Arjovsky et al.\ \cite{arjovsky2019irm}
& Out-of-distribution generalisation
& IRM: invariant risk minimisation
& Multiple labelled environments
& Encourages causally stable features. Guarantees weaken when the environment count is small. \\
\hline
Perry et al.\ \cite{perry2022causal}
& Multi-environment causal graph recovery
& Sparse Mechanism Shift hypothesis with MSS score
& Multiple labelled environments
& Full graph identifiable under sparse shifts. Environment membership must still be known. \\
\hline
Huang et al.\ \cite{huang2019specific}
& Causal clustering with shared and specific relations
& SSCM: mixture of causal mechanisms
& i.i.d., multiple regimes
& Distinguishes shared and regime-specific causal relations. Requires a pre-specified regime count. \\
\hline
Chen et al.\ \cite{chen2023ccsl}
& Causal structure learning from unknown environments
& CCSL: Bayesian nonparametric Causal CRP with variational inference
& i.i.d., unknown environments
& Does not require a regime count in advance. Identifiable under a linear non-Gaussian assumption. Primary baseline for this thesis. \\
\hline
Varambally et al.\ \cite{varambally2024mcd}
& Mixture of causal models from time series
& MCD: variational inference over mixtures of SCMs
& Temporal, autoregressive structure
& Jointly infers causal models and mixing probabilities. Relies on temporal ordering. \\
\hline
Rahmani and Frossard \cite{rahmani2025castor}
& Temporal causal regime discovery
& CASTOR: EM over piecewise-constant regimes with NOTEARS acyclicity
& Sequential time series
& Jointly learns regime boundaries and per-regime DAGs without a prior regime count. Relies on temporal contiguity unavailable in i.i.d.\ settings. \\
\hline
Cheng et al.\ \cite{cheng2025dycast}
& Continuously evolving causal structure
& DyCAST: Neural ODE constrained to DAG manifold
& Sequential time series
& Captures smooth causal transitions and scales beyond 50 nodes. Requires temporal ordering. \\
\hline
Li et al.\ \cite{li2025hcl}
& Unsupervised causal clustering from mixed-type data
& HCL: bi-directional iterative causal clustering (preprint)
& i.i.d., no environment labels
& No temporal ordering or regime count required. Scalability on large graphs has not yet been assessed. \\
\hline
Sadeghi et al.\ \cite{sadeghi2025causal}
& Causal discovery from nonstationary time series
& CD-NOTS: nonparametric extension of CD-NOD correcting lagged conditioning
& Non-stationary time series
& Uncovers nonlinear nonstationary causal relations missed by stationarity-assuming methods. \\
\hline
Mooij et al.\ \cite{mooij2020jci}
& Joint causal inference from multiple contexts
& JCI: context encoded as additional graph nodes
& Multiple pre-specified contexts
& Theoretically general and subsumes many special cases. Contexts must be labelled in advance. \\
\hline
Faltenbacher et al.\ \cite{faltenbacher2025internal}
& Assumption violation detection
& Internal incoherency scores for constraint-based algorithms
& i.i.d., single environment
& Flags assumption violations without ground truth. Not yet adapted to causal clustering settings. \\
\hline

\end{longtable}
}

